\documentclass[12pt]{article}
\usepackage{amsmath}
\usepackage{amssymb}
\usepackage{physics}
\usepackage{hyperref}

% Shorthand for displaystyle fractions inside inline math
\newcommand{\ds}{\displaystyle}

\title{Composition of Velocities}
\author{}
\date{}

\begin{document}
\maketitle

% ============================================================
\section{Introduction}
% ============================================================

The preceding essay \href{/physics/relativity/special-relativity/lorentz-transformations}{Lorentz Transformations}
derives the transformation between two inertial reference frames $F$ and $F'$, where $F'$ moves
with velocity $\vb{v}=v\vb{e}_1$ relative to $F$. For the derivation below, we need the Lorentz
transformation written explicitly in terms of this relative velocity:

\begin{equation}
  \boldsymbol{\Lambda}(\vb{v})
  =
    \frac{\gamma}{v^2}\ket{\vb{v},0}\bra{\vb{v},0}
    - \frac{\gamma}{c}\ket{\vb{v},0}\bra{\vb{0},1}
    + \ket{\vb{e}_2,0}\bra{\vb{e}_2,0}
    + \ket{\vb{e}_3,0}\bra{\vb{e}_3,0}
    - \frac{\gamma}{c}\ket{\vb{0},1}\bra{\vb{v},0}
    + \gamma\ket{\vb{0},1}\bra{\vb{0},1},
  \label{eq:Lorentz_transformation_in_velocity_form}
\end{equation}

where

\[
  \gamma
  =
  \frac{1}{\sqrt{1-\dfrac{v^2}{c^2}}}.
\]

The time-coordinate component of this transformation is

\begin{equation}
  ct'
  =
  \gamma\!\left(
    -\frac{v}{c}\ip{\vb{e}_1}{\vb{x}} + ct
  \right).
  \label{eq:Lorentz_time_coordinate}
\end{equation}

% ============================================================
\section{Derivation of Composition of Velocities}
% ============================================================

In classical mechanics, a velocity $\vb{u}'$ measured in the inertial frame $F'$ is measured
in the inertial frame $F$ as

\[
\vb{u} = \vb{v} + \vb{u}',
\]

which is the familiar Galilean law of velocity addition. In relativity, however, velocities
must be composed according to the Lorentz transformation, with the Galilean law recovered
in the limit in which the relevant velocities are much smaller than the speed of light.

We consider the inertial reference frames $F$ and $F'$ defined in the introduction and the
general motion of a body in $F'$ described by the spacetime trajectory

\[
  \ket{\vb{x}',0} + ct'\ket{\vb{0},1},
\]

where we assume that $\vb{x}' = \vb{x}'(t')$. The motion of the body observed in $F$ is
related to the motion observed in $F'$ through the Lorentz transformation

\[
  \boldsymbol{\Lambda}(\vb{v})
  \left(
    \ket{\vb{x},0} + ct\ket{\vb{0},1}
  \right)
  =
  \ket{\vb{x}',0} + ct'\ket{\vb{0},1},
\]

where $\boldsymbol{\Lambda}(\vb{v})$ is given by
\eqref{eq:Lorentz_transformation_in_velocity_form} and $\vb{x}=\vb{x}(t)$.
Differentiating the spacetime trajectory with respect to the
time $t$ measured in $F$ yields

\[
  \boldsymbol{\Lambda}(\vb{v})\,
  \frac{d}{dt}
  \left(
    \ket{\vb{x},0} + ct\ket{\vb{0},1}
  \right)
  =
  \frac{d}{dt}
  \left(
    \ket{\vb{x}',0} + ct'\ket{\vb{0},1}
  \right).
\]

%%%%%%%%%%%%%%%%%%%%%%%%%%%%%%%%%%%%%%%

Differentiating the time-coordinate transformation
\eqref{eq:Lorentz_time_coordinate}, we obtain

\begin{equation}
  \frac{dt'}{dt}
  =
  \gamma\!\left(
    1-\frac{\ip{\vb{v}}{\vb{u}}}{c^2}
  \right).
  \label{eq:time_derivative_velocity_composition}
\end{equation}

Since

\[
  \frac{d\vb{x}'}{dt}
  =
  \frac{d\vb{x}'}{dt'}\frac{dt'}{dt}
  =
  \vb{u}'\frac{dt'}{dt},
\]

the differentiated Lorentz transformation becomes

\[
  \boldsymbol{\Lambda}(\vb{v})
  \bigl(
    \ket{\vb{u},0}+c\ket{\vb{0},1}
  \bigr)
  =
  \gamma\!\left(
    1-\frac{\ip{\vb{v}}{\vb{u}}}{c^2}
  \right)
  \bigl(
    \ket{\vb{u}',0}+c\ket{\vb{0},1}
  \bigr).
\]

Substituting \eqref{eq:Lorentz_transformation_in_velocity_form}, we obtain

\[
  \begin{aligned}
    &\gamma\!\left(
      \frac{\ip{\vb{v}}{\vb{u}}}{v^2} - 1
    \right)\ket{\vb{v},0}
    + \ket{\vb{e}_2,0}\ip{\vb{e}_2}{\vb{u}}
    + \ket{\vb{e}_3,0}\ip{\vb{e}_3}{\vb{u}}
    + \ket{\vb{0},1}\,
      \gamma c\!\left(
        1 - \frac{\ip{\vb{v}}{\vb{u}}}{c^2}
      \right)
    \\
    &=
    \gamma\!\left(
      1 - \frac{\ip{\vb{v}}{\vb{u}}}{c^2}
    \right)
    \bigl(
      \ket{\vb{u}',0}+c\ket{\vb{0},1}
    \bigr).
  \end{aligned}
\]

The time components on both sides are identical and therefore cancel, leaving

\[
  \gamma\!\left(
    \frac{\ip{\vb{v}}{\vb{u}}}{v^2} - 1
  \right)\ket{\vb{v},0}
  + \ket{\vb{e}_2,0}\ip{\vb{e}_2}{\vb{u}}
  + \ket{\vb{e}_3,0}\ip{\vb{e}_3}{\vb{u}}
  =
  \gamma\!\left(
    1 - \frac{\ip{\vb{v}}{\vb{u}}}{c^2}
  \right)\ket{\vb{u}',0}.
\]

Using the decomposition

\[
  \ket{\vb{u},0}
  =
  \frac{\ip{\vb{v}}{\vb{u}}}{v^2}\ket{\vb{v},0}
  +
  \ket{\vb{e}_2,0}\ip{\vb{e}_2}{\vb{u}}
  +
  \ket{\vb{e}_3,0}\ip{\vb{e}_3}{\vb{u}},
\]

the previous equation can be written more compactly as

\[
  \gamma\,
  \frac{\ip{\vb{v}}{\vb{u}}}{v^2}\ket{\vb{v},0}
  -
  \gamma\ket{\vb{v},0}
  -
  \frac{\ip{\vb{v}}{\vb{u}}}{v^2}\ket{\vb{v},0}
  +
  \ket{\vb{u},0}
  =
  \gamma\!\left(
    1 - \frac{\ip{\vb{v}}{\vb{u}}}{c^2}
  \right)\ket{\vb{u}',0},
\]

or

\begin{equation}
  (\gamma - 1)
  \frac{\ip{\vb{v}}{\vb{u}}}{v^2}\ket{\vb{v},0}
  -
  \gamma\ket{\vb{v},0}
  +
  \ket{\vb{u},0}
  =
  \gamma\!\left(
    1 - \frac{\ip{\vb{v}}{\vb{u}}}{c^2}
  \right)\ket{\vb{u}',0}.
  \label{eq:velocities_relationship}
\end{equation}

Solving for $\ket{\vb{u}',0}$ gives

\begin{equation}
  \ket{\vb{u}',0}
  =
  \frac{1}{
    \displaystyle
    1 - \frac{\ip{\vb{v}}{\vb{u}}}{c^2}
  }
  \left(
    \frac{\gamma-1}{\gamma}\,
    \frac{\ip{\vb{v}}{\vb{u}}}{v^2}\ket{\vb{v},0}
    -
    \ket{\vb{v},0}
    +
    \frac{\ket{\vb{u},0}}{\gamma}
  \right).
  \label{eq:inverse_velocity_composition}
\end{equation}

%%%%%%%%%%%%%%%%%%%%%%%%%%%%%%%%%%%%%%%%%%%%%%%%%%%%%%%%%%%%%%%

The velocity $\ket{\vb{u}}$ can be easily expressed in terms of the velocity $\ket{\vb{u}'}$
by exchanging the roles of $\ket{\vb{u}}$ and $\ket{\vb{u}'}$ in the equation above, and
changing the sign of the relative velocity $\vb{v}$.

As an exercise, we derive the same result by directly inverting
\eqref{eq:velocities_relationship}.

We start by factoring the velocity $\ket{\vb{u}}$ in \eqref{eq:velocities_relationship}
and express the equation as an operator on the velocity $\ket{\vb{u}}$:

\[
  \left(
    \frac{\gamma}{v^2}\ket{\vb{v}}\bra{\vb{v}}
    + \ket{\vb{e}_2,0}\bra{\vb{e}_2,0}
    + \ket{\vb{e}_3,0}\bra{\vb{e}_3,0}
    + \frac{\gamma}{c^2}\ket{\vb{u}'}\bra{\vb{v}}
  \right)\ket{\vb{u}}
  =
  \gamma\bigl(\ket{\vb{u}'} + \ket{\vb{v}}\bigr).
\]

Expressing $\ket{\vb{u}'}$ in the basis $\{\vb{e}_1,\vb{e}_2,\vb{e}_3\}$, the operator becomes

\[
  \begin{aligned}
    \Bigl(
      &\gamma\!\left(1 + \frac{v}{c^2}u'_1\right)
        \ket{\vb{e}_1,0}\bra{\vb{e}_1,0}
      + \gamma\,\frac{v}{c^2}u'_2\,
        \ket{\vb{e}_2,0}\bra{\vb{e}_1,0}
      + \gamma\,\frac{v}{c^2}u'_3\,
        \ket{\vb{e}_3,0}\bra{\vb{e}_1,0}
      \\
      &+ \ket{\vb{e}_2,0}\bra{\vb{e}_2,0}
      + \ket{\vb{e}_3,0}\bra{\vb{e}_3,0}
    \Bigr)\cdot\ket{\vb{u},0}
    =
    \gamma\bigl(\ket{\vb{u}'} + \ket{\vb{v}}\bigr).
  \end{aligned}
\]

The operator

\[
  \gamma\!\left(1 + \frac{v}{c^2}u'_1\right)
  \ket{\vb{e}_1,0}\bra{\vb{e}_1,0}
  + \gamma\,\frac{v}{c^2}u'_2\,
    \ket{\vb{e}_2,0}\bra{\vb{e}_1,0}
  + \gamma\,\frac{v}{c^2}u'_3\,
    \ket{\vb{e}_3,0}\bra{\vb{e}_1,0}
  + \ket{\vb{e}_2,0}\bra{\vb{e}_2,0}
  + \ket{\vb{e}_3,0}\bra{\vb{e}_3,0}
\]

has the inverse

\[
  \begin{aligned}
    &\frac{1}{
      \displaystyle
      \gamma\!\left(1 + \frac{v}{c^2}u'_1\right)}
      \ket{\vb{e}_1,0}\bra{\vb{e}_1,0}
    - \frac{v}{c^2}
      \frac{u'_2}{
        \displaystyle
        1 + \frac{v}{c^2}u'_1}
      \ket{\vb{e}_2,0}\bra{\vb{e}_1,0}
    \\
    &- \frac{v}{c^2}
      \frac{u'_3}{
        \displaystyle
        1 + \frac{v}{c^2}u'_1}
      \ket{\vb{e}_3,0}\bra{\vb{e}_1,0}
    + \ket{\vb{e}_2,0}\bra{\vb{e}_2,0}
    + \ket{\vb{e}_3,0}\bra{\vb{e}_3,0}.
  \end{aligned}
\]

We can express the velocity $\ket{\vb{u}}$ as

\[
  \begin{aligned}
    \ket{\vb{u},0} =\;
    &\left(
      \frac{1}{
        \displaystyle
        \gamma\!\left(1 + \frac{v}{c^2}u'_1\right)}
        \ket{\vb{e}_1,0}\bra{\vb{e}_1,0}
      - \frac{v}{c^2}
        \frac{u'_2}{
          \displaystyle
          1 + \frac{v}{c^2}u'_1}
        \ket{\vb{e}_2,0}\bra{\vb{e}_1,0}
      \right.
      \\
    &\left.
      - \frac{v}{c^2}
        \frac{u'_3}{
          \displaystyle
          1 + \frac{v}{c^2}u'_1}
        \ket{\vb{e}_3,0}\bra{\vb{e}_1,0}
      + \ket{\vb{e}_2,0}\bra{\vb{e}_2,0}
      + \ket{\vb{e}_3,0}\bra{\vb{e}_3,0}
    \right)
    \\[0.4em]
    &\cdot\left(
      v\gamma\ket{\vb{e}_1,0}
      + \gamma\ket{\vb{u}'}
    \right),
  \end{aligned}
\]

which simplifies to

\[
  \ket{\vb{u},0}
  =
  \frac{1}{
    \displaystyle
    1 + \frac{v}{c^2}u'_1}
  \left(
    (v + u'_1)\ket{\vb{e}_1,0}
    + \sqrt{1 - \frac{v^2}{c^2}}\,
      u'_2\,\ket{\vb{e}_2,0}
    + \sqrt{1 - \frac{v^2}{c^2}}\,
      u'_3\,\ket{\vb{e}_3,0}
  \right).
\]

Suppressing the zero time component in the notation, we can write the previous equation
in the more compact vector form

\begin{equation}
  \ket{\vb{u}}
  =
  \frac{1}{
    \displaystyle
    1 + \frac{\ip{\vb{v}}{\vb{u}'}}{c^2}}
  \left(
    \frac{\gamma - 1}{\gamma}\,
    \frac{\ip{\vb{u}'}{\vb{v}}}{v^2}\ket{\vb{v}}
    + \ket{\vb{v}}
    + \frac{\ket{\vb{u}'}}{\gamma}
  \right).
  \label{eq:velocity_composition}
\end{equation}

As a final remark, we should note two important cases. First, from
equation \eqref{eq:velocity_composition}, we notice that the velocity of any
object at rest in $F'$ is observed in $F$ to move with velocity $\ket{\vb{v}}$,
as expected. Conversely, any object at rest in $F$ is observed in $F'$ to move
with velocity $-\ket{\vb{v}}$, as follows from
equation \eqref{eq:inverse_velocity_composition}.

The important point is that, although the velocity vector changes sign when the roles of
 the two inertial frames are exchanged, the magnitude of their relative velocity does not.
 If $F'$ moves with velocity $\vb{v}$ relative to $F$, then $F$ moves with velocity $-\vb{v}$ relative to $F'$,
  and therefore both frames assign the same relative speed $v$ to their motion. Consequently, they also
  associate the same Lorentz factor $\gamma(v)$ with their relative motion. This reciprocity will play an
  important role in the discussion of the twin paradox.

% ============================================================
\section{Limiting Cases}
% ============================================================

In this section we derive the limiting formulas for the composition of velocities when one
of the relevant velocities approaches the speed of light.

\subsection{A Particle Moving at the Speed of Light in $F$}

We start by considering the case where $\ket{\vb{u}}$ approaches the speed of light
$c$ and determine the value that $\ket{\vb{u}'}$ reaches. To achieve this, we compute
the square modulus of $\ket{\vb{u}'}$ from \eqref{eq:inverse_velocity_composition},
which reads

\[
  u'^2
  =
  \frac{1}{
    \displaystyle
    \left(
      1 - \frac{\ip{\vb{v}}{\vb{u}}}{c^2}
    \right)^2
  }
  \left(
    \frac{\ip{\vb{u}}{\vb{v}}^2}{c^2}
    + v^2
    + \frac{u^2}{\gamma^2}
    - 2\ip{\vb{u}}{\vb{v}}
  \right).
\]

When $\ket{\vb{u}} = c\ket{\vb{n}}$, where $\ket{\vb{n}}$ is a unit vector, the equation
above simplifies to

\[
  u'^2
  =
  \frac{c^2}{
    \displaystyle
    \left(
      1 - \frac{\ip{\vb{v}}{\vb{n}}}{c}
    \right)^2
  }
  \left(
    1 - \frac{\ip{\vb{v}}{\vb{n}}}{c}
  \right)^2
  =
  c^2,
\]

which implies that the velocity $\ket{\vb{u}'}$ reaches the speed of light when
$\ket{\vb{u}}$ approaches the speed of light.

\subsection{A Particle Moving at the Speed of Light in $F'$}

Similarly, we can determine the limit of the velocity $\ket{\vb{u}}$ when the velocity
$\ket{\vb{u}'}$ approaches the speed of light. From the velocity composition formula
\eqref{eq:velocity_composition}, the modulus of the velocity $\ket{\vb{u}}$ reads

\[
  u^2
  =
  \frac{1}{
    \displaystyle
    \left(
      1 + \frac{\ip{\vb{v}}{\vb{u}'}}{c^2}
    \right)^2
  }
  \left(
    \frac{\ip{\vb{u}'}{\vb{v}}^2}{c^2}
    + v^2
    + \frac{u'^2}{\gamma^2}
    + 2\ip{\vb{u}'}{\vb{v}}
  \right).
\]

This differs from the previous expression only by the replacement
$\vb{v}\rightarrow-\vb{v}$. When $\ket{\vb{u}'} = c\ket{\vb{n}}$, where
$\ket{\vb{n}}$ is a unit vector, the equation above simplifies to

\[
  u^2
  =
  \frac{c^2}{
    \displaystyle
    \left(
      1 + \frac{\ip{\vb{v}}{\vb{n}}}{c}
    \right)^2
  }
  \left(
    1 + \frac{\ip{\vb{v}}{\vb{n}}}{c}
  \right)^2
  =
  c^2,
\]

which, also in this case, means that the velocity $\ket{\vb{u}}$ reaches the speed of light
when $\ket{\vb{u}'}$ approaches the speed of light.

\subsection{Relative Motion Approaching the Speed of Light}

Finally, consider the limit in which the magnitude of the relative velocity approaches the
speed of light,

\[
  \ket{\vb{v}} \longrightarrow c\ket{\vb{n}},
\]

where $\ket{\vb{n}}$ is a unit vector. In this limit, the modulus of the velocity
$\ket{\vb{u}}$ becomes

\[
  u^2
  \longrightarrow
  \frac{c^2}{
    \displaystyle
    \left(
      1 + \frac{\ip{\vb{u}'}{\vb{n}}}{c}
    \right)^2
  }
  \left(
    1 + \frac{\ip{\vb{u}'}{\vb{n}}}{c}
  \right)^2
  =
  c^2,
\]

independently of the value of $\ket{\vb{u}'}$.

\end{document}
